\documentclass[12pt]{article}
\usepackage[utf8]{inputenc}
\usepackage[T1]{fontenc}
\usepackage[margin=1in]{geometry}
\usepackage{setspace}
\usepackage{xcolor}
\usepackage[colorlinks=true, urlcolor=blue!60!black, linkcolor=blue!60!black]{hyperref}
\setlength{\parindent}{2em}
\setlength{\parskip}{1.25em}
\renewcommand{\baselinestretch}{1.0}
\title{\textit{Baker v. General Motors Corp.}, 522 U.S. 222 (1998)}
\author{Case Brief}
\date{}
\begin{document}
\maketitle

\section{Case Metadata}
\begin{itemize}
  \item \textbf{Case Name:} \textit{Baker v. General Motors Corp.}, 522 U.S. 222 (1998) (No.~96-653)
  \item \textbf{Chapter:} Chapter 4: Federalism and State Sovereignty
  \item \textbf{Court and Dates:} Supreme Court of the United States; argued October 15, 1997; decided January 13, 1998
  \item \textbf{Lower Court:} U.S. Court of Appeals for the Eighth Circuit, 86 F.3d 811 (8th Cir. 1996) (reversed and remanded)
  \item \textbf{Vote:} 9--0 in the judgment (5--4 split on doctrinal rationale)
  \item \textbf{Majority Author:} Justice Ruth Bader Ginsburg (joined by Chief Justice Rehnquist and Justices Stevens, Souter, and Breyer)
  \item \textbf{Concurrences in the Judgment:} Justice Antonin Scalia; Justice Anthony M. Kennedy (joined by Justices O'Connor and Thomas)
  \item \textbf{Primary Texts:} \href{https://www.law.cornell.edu/supremecourt/text/96-653}{Cornell LII Full Opinion}; \href{https://supreme.justia.com/cases/federal/us/522/222/}{Justia U.S. Supreme Court Center}
  \item \textbf{Constitutional and Statutory Provisions:} U.S. Const. art.~IV, \S~1 (Full Faith and Credit Clause); Full Faith and Credit Act, 28 U.S.C. \S~1738
\end{itemize}

\section{Facts and Procedural History}
Two lawsuits in different states generated the Full Faith and Credit dispute in \href{https://www.law.cornell.edu/supct/html/96-653.ZO.html}{\textit{Baker v. General Motors Corp.}} For fifteen years, Ronald Elwell worked as an engineering analyst for General Motors Corporation (GM), studying vehicular fires and frequently assisting GM counsel in defending product liability suits. After Elwell's employment relationship with GM soured, he sued GM in a Michigan county court for wrongful discharge, and GM counterclaimed. On August 26, 1992, the parties settled: GM paid Elwell an undisclosed sum, and the Michigan court entered a stipulated permanent injunction barring Elwell from ``testifying, without the prior written consent of [GM], \ldots as \ldots a witness of any kind \ldots in any litigation already filed, or to be filed in the future, involving [GM] as an owner, seller, manufacturer and/or designer \ldots'' (\href{https://www.law.cornell.edu/supct/html/96-653.ZO.html}{522 U.S. at 228}). The consent decree carved out only a then-pending Georgia fuel-tank action, though a separate GM--Elwell settlement agreement stipulated that if another court or tribunal ordered Elwell to testify, his testimony would ``in no way'' form a basis for a contempt or breach-of-contract action by GM (\href{https://www.law.cornell.edu/supct/html/96-653.ZO.html}{id. at 228--29}).

Meanwhile, Beverly Garner died in a February 1990 Missouri highway crash and vehicle fire while riding as a front-seat passenger in a 1985 Chevrolet S-10 Blazer. Her sons, Kenneth Lee Baker and Steven Robert Baker (appearing by his next friend, Melissa Thomas), filed a wrongful-death product liability suit against GM in Missouri state court, alleging a faulty GM fuel pump; GM removed the action to federal district court (\href{https://www.law.cornell.edu/supremecourt/text/96-653}{522 U.S. at 229--30}). The Bakers, who had no connection to the Michigan litigation, subpoenaed Elwell. After \textit{in camera} review, the federal district court permitted Elwell's deposition and trial testimony on two alternative grounds: (1) enforcing Michigan's gag injunction would violate Missouri's public policy favoring disclosure of relevant, nonprivileged evidence, and (2) the injunction was modifiable in Missouri just as in Michigan. Elwell testified, and a jury awarded the Bakers \$11.3 million (\href{https://www.law.cornell.edu/supremecourt/text/96-653}{id. at 230}). The Eighth Circuit reversed (86 F.3d 811), concluding that even assuming a public-policy exception existed, Missouri's interest in Full Faith and Credit matched its disclosure policy, and the record did not show Michigan would modify its decree. The Supreme Court granted certiorari.

\section{Issue(s) Presented}
Whether the Full Faith and Credit Clause (U.S. Const. art.~IV, \S~1) and its implementing statute (28 U.S.C. \S~1738) bar nonparty plaintiffs in a Missouri federal wrongful-death action from subpoenaing a witness whose testimony against the defendant is prohibited by a stipulated permanent injunction entered by a Michigan state court.

\section{Holding}
No. The Court reversed and remanded unanimously in the judgment (9--0), holding that Ronald Elwell could testify in the Missouri action without offending the Full Faith and Credit Clause or 28 U.S.C. \S~1738. Writing for a five-Justice majority, Justice Ginsburg reaffirmed that the Full Faith and Credit obligation for final judgments is exacting and admits of no roving ``public policy exception,'' sharply distinguishing judgments from choice-of-law rules for statutes. Nevertheless, a sister-state equitable decree cannot dictate evidentiary rules or override another jurisdiction's subpoena power in litigation brought by strangers over whom the rendering court lacked adjudicatory authority (\href{https://www.law.cornell.edu/supremecourt/text/96-653}{522 U.S. at 233--41}).

\section{Reasoning}
Justice Ginsburg structured the Court's opinion around four doctrinal pillars grounded in Article~IV, Section~1 and the 1790 Full Faith and Credit Act:

\begin{enumerate}
  \item \textbf{Sharp Distinction Between Laws and Judgments:} The Clause's animating purpose is to transform independent states into ``integral parts of a single nation'' (\textit{Milwaukee County v. M. E. White Co.}, 296 U.S. 268, 277 (1935)). For \textit{laws} (statutes and common law), Full Faith and Credit is relatively permissive: a forum state need not substitute a sister state's statutes for its own on a subject over which it may legislate, and a forum court may consult local ``public policy'' in choice-of-law analysis (\href{https://www.law.cornell.edu/supremecourt/text/96-653}{522 U.S. at 232--33}). For \textit{final judgments}, however, the command is ``exacting'': once rendered by a court with subject-matter and personal jurisdiction, a judgment commands nationwide claim and issue preclusion, with \textit{no roving public-policy exception}---even where the underlying claim would be unenforceable in the forum state (as in \textit{Fauntleroy v. Lum}, 210 U.S. 230 (1908)).
  \item \textbf{Equity Decrees Are Not Categorically Exempt, but Enforcement Mechanisms Do Not Travel:} Rejecting any blanket carve-out for equitable decrees, the Court noted that money decrees in equity long enjoyed nationwide recognition. Yet Full Faith and Credit does not compel a sister state to adopt the rendering state's ``time, manner, and mechanisms for enforcing judgments'' (\href{https://www.law.cornell.edu/supremecourt/text/96-653}{522 U.S. at 235}). Enforcement measures remain under the evenhanded control of forum law.
  \item \textbf{Territorial and Jurisdictional Limits on Injunctive Decrees:} Sister-state orders commanding action or inaction are denied extraterritorial effect when they purport to accomplish an official act within another state's exclusive province (such as directly transferring title to foreign land, \textit{Fall v. Eastin}, 215 U.S. 1 (1909)) or interfere with litigation over which the rendering court lacked authority. While the Michigan decree validly barred Elwell from \textit{volunteering} as an expert against GM and bound GM and Elwell in Michigan, Michigan had no jurisdiction over the Bakers and no authority to force them to intervene in Michigan (\textit{Martin v. Wilks}, 490 U.S. 755 (1989)).
  \item \textbf{Control Over Evidentiary Subpoenas and Nonparty Litigation:} Controlling witness competence, relevance, and subpoena power in suits brought by strangers to the Michigan proceeding lies outside Michigan's governance (\href{https://supreme.justia.com/cases/federal/us/522/222/}{522 U.S. at 238--39}). The consent decree itself acknowledged this limit by carving out the pending Georgia suit, and GM's side agreement promising not to pursue Michigan contempt against Elwell if subpoenaed elsewhere underscored that GM sought to turn Missouri into a ``surrogate tribunal'' to enforce an injunction Michigan lacked power to impose on nonparties (\href{https://www.law.cornell.edu/supct/html/96-653.ZO.html}{id. at 240--41}).
\end{enumerate}

\section{Concurring and Dissenting Opinions}
No Justice dissented, but four Justices rejected the majority's doctrinal path in two concurrences in the judgment:

\begin{itemize}
  \item \textbf{Justice Scalia (concurring in the judgment):} Writing alone, \href{https://www.law.cornell.edu/supct/html/96-653.ZC.html}{Justice Scalia} urged a far simpler, historical rule drawn from \textit{McElmoyle ex rel. Bailey v. Cohen}, 38 U.S. (13 Pet.) 312 (1839), and \textit{Thompson v. Whitman}, 85 U.S. (18 Wall.) 457 (1874). Full Faith and Credit establishes a rule of evidence and preclusion, not extraterritorial execution: a state court's judgment ``cannot be enforced out of the state by an execution issued within it.'' Because GM asked the Missouri federal court directly to execute Michigan's injunction by silencing Elwell, Missouri had no constitutional obligation to do so.
  \item \textbf{Justice Kennedy (concurring in the judgment, joined by Justices O'Connor and Thomas):} \href{https://www.law.cornell.edu/supct/html/96-653.ZC1.html}{Justice Kennedy} warned that the majority's two broad exceptions---for decrees affecting exclusive sister-state official acts or interfering with sister-state litigation---sat in ``uneasy tension'' with its disavowal of a public-policy exception and needlessly unsettled equity and antisuit-injunction doctrine. Under \textit{Kremer v. Chemical Construction Corp.}, 456 U.S. 461 (1982), and \textit{Matsushita Electric Industrial Co. v. Epstein}, 516 U.S. 367 (1996), the threshold inquiry under 28 U.S.C. \S~1738 is the preclusive effect the judgment has in the \textit{rendering} state. Because Michigan law requires mutuality of estoppel and does not bind nonparties to prior injunctions, Michigan itself would not preclude the Bakers; that state-law answer resolved the case without creating new constitutional exceptions.
\end{itemize}

\section{Key Quotations}
\begin{quote}
``Regarding judgments, however, the full faith and credit obligation is exacting. A final judgment in one State, if rendered by a court with adjudicatory authority over the subject matter and persons governed by the judgment, qualifies for recognition throughout the land. \ldots But our decisions support no roving `public policy exception' to the full faith and credit due judgments.'' --- Justice Ginsburg, \href{https://www.law.cornell.edu/supremecourt/text/96-653}{522 U.S. at 233}
\end{quote}

\begin{quote}
``Full faith and credit, however, does not mean that States must adopt the practices of other States regarding the time, manner, and mechanisms for enforcing judgments. Enforcement measures do not travel with the sister state judgment as preclusive effects do; such measures remain subject to the evenhanded control of forum law.'' --- Justice Ginsburg, \href{https://www.law.cornell.edu/supremecourt/text/96-653}{522 U.S. at 235}
\end{quote}

\begin{quote}
``Most essentially, Michigan lacks authority to control courts elsewhere by precluding them, in actions brought by strangers to the Michigan litigation, from determining for themselves what witnesses are competent to testify and what evidence is relevant and admissible in their search for the truth.'' --- Justice Ginsburg, \href{https://supreme.justia.com/cases/federal/us/522/222/}{522 U.S. at 238}
\end{quote}

\begin{quote}
``The Court's reliance upon unidentified principles to justify omitting certain types of injunctions from the doctrine's application leaves its decision in uneasy tension with its own rejection of a broad public policy exception to full faith and credit.'' --- Justice Kennedy (concurring in the judgment), \href{https://www.law.cornell.edu/supct/html/96-653.ZC1.html}{522 U.S. at 245}
\end{quote}

\section{Doctrinal Significance and Modern Legacy (Through 2026)}
\textit{Baker} remains the Supreme Court's canonical synthesis of Article~IV, Section~1, anchoring two enduring rules: (1) the sharp asymmetry between permissive choice-of-law review of sister-state statutes and the ``exacting'' recognition mandate for sister-state judgments, and (2) the boundary between binding \textit{res judicata} and non-traveling \textit{enforcement mechanisms}.

\begin{itemize}
  \item \textbf{Post-\textit{Baker} Supreme Court Reaffirmations:} In \href{https://supreme.justia.com/cases/federal/us/538/488/case.pdf}{\textit{Franchise Tax Board of California v. Hyatt} (\textit{Hyatt I}), 538 U.S. 488 (2003)}, the Court relied on \textit{Baker}'s law-versus-judgment divide to hold that Nevada did not violate Full Faith and Credit by declining to apply California's sovereign-immunity statute (though \href{https://supreme.justia.com/cases/federal/us/587/17-1299/}{\textit{Hyatt III}, 587 U.S. 230 (2019)}, later overruled \textit{Nevada v. Hall} on structural interstate sovereign-immunity grounds rather than Full Faith and Credit). In \href{https://www.law.cornell.edu/supremecourt/text/05-555}{\textit{Lance v. Dennis}, 546 U.S. 459 (2006)}, the Court cited \textit{Baker} to keep nonparty preclusion governed by 28 U.S.C. \S~1738 rather than an expanded \textit{Rooker-Feldman} bar. Most decisively, in \href{https://supreme.justia.com/cases/federal/us/577/464/}{\textit{V.L. v. E.L.}, 577 U.S. 464 (2016) (per curiam)}, the Court invoked \textit{Baker}'s ``exacting'' judgment rule to summarily reverse the Alabama Supreme Court's refusal to recognize a Georgia same-sex second-parent adoption decree.
  \item \textbf{Records, Executive Officials, and Congress's Effects Power:} Lower courts have wrestled with \textit{Baker}'s enforcement-mechanism boundary. In \textit{Adar v. Smith}, 639 F.3d 146 (5th Cir. 2011) (en banc), the Fifth Circuit cited \textit{Baker} in holding that Louisiana's executive registrar need not issue a revised birth certificate reflecting a New York adoption decree (\href{https://constitutioncenter.org/the-constitution/articles/article-iv/clauses/44}{National Constitution Center}). Meanwhile, Congress exercised its Article~IV, Section~1 ``Effects'' power in the Respect for Marriage Act of 2022 (28 U.S.C. \S~1738C) to mandate interstate recognition of valid same-sex and interracial marriages.
  \item \textbf{Interstate Shield-Law Litigation (2025--2026):} Following \textit{Dobbs}, \textit{Baker} sits at the center of high-stakes interstate judgment-enforcement clashes. After a Texas court entered a February 2025 default judgment exceeding \$113,000 and an injunction against New York physician Dr.~Margaret Carpenter for telehealth abortion prescriptions, Ulster County Clerk Taylor Bruck refused to file the judgment under New York's 2023 shield law; in late 2025, a New York court dismissed Texas's mandamus action (\href{https://journals.law.umn.edu/mjlst/2026/11/10/telehealth-abortion-and-the-limits-of-state-shield-laws}{Minn. J.L. Sci. \& Tech., Nov. 2026}). Through 2026, courts and scholars actively debate whether \textit{Baker}'s no-public-policy rule compels enforcement of such out-of-state money judgments, or whether the penal-judgment exception, lack of personal jurisdiction, or \textit{Baker}'s limits on extraterritorial injunctions and forum enforcement mechanisms shield sanctuary states.
\end{itemize}

\section{Pedagogical Critiques and Exam Hypotheticals}
\textbf{Pedagogical Critique:} Justice Kennedy's concurrence exposes the core analytical vulnerability in Justice Ginsburg's opinion: if Full Faith and Credit truly admits of ``no roving public policy exception'' for judgments, carving out injunctions that ``interfere with litigation'' in a sister state risks smuggling a functional public-policy exception in through a jurisdictional label. Conversely, Justice Ginsburg's Footnote~12 (\href{https://www.law.cornell.edu/supct/html/96-653.ZO.html}{522 U.S. at 239 n.12}) highlights the flaw in Justice Kennedy's pure-preclusion approach: GM argued it was only enforcing an injunction against \textit{Elwell} (a Michigan party), not precluding the Bakers---meaning state mutuality doctrine alone could not cleanly explain why a Missouri court may compel Elwell to defy a valid Michigan order.

\textbf{Exam Hypothetical (Standard Application):} State~A enters a stipulated permanent injunction barring a former pharmaceutical scientist from testifying nationwide about a drug's side effects, with a side letter waiving contempt if subpoenaed. Nonparty plaintiffs injured in State~B sue the manufacturer in State~B and subpoena the scientist. \textit{Analysis:} Under \textit{Baker}, State~B's court may compel testimony without violating Article~IV, Section~1; State~A lacks authority to control evidentiary subpoenas in nonparty litigation outside its borders.

\textbf{``Opposite Side'' Contrast (When Full Faith and Credit \textit{Does} Bind):} Suppose instead that the Bakers had intervened in the Michigan proceeding, fully litigated whether Elwell's testimony was protected by attorney--client privilege, and lost in a final Michigan judgment. Or suppose a State~A court with personal and subject-matter jurisdiction enters a final \textit{money judgment} for breach of a non-disclosure settlement against Elwell, and GM seeks to domesticate that money judgment in State~B, whose legislature has declared such confidentiality clauses contrary to State~B's public policy. \textit{Contrast:} In both scenarios, \textit{Baker}'s ``Opposite Side'' rule controls---as Justice Ginsburg emphasized in Footnote~12 and her discussion of \textit{Fauntleroy v. Lum} (\href{https://www.law.cornell.edu/supremecourt/text/96-653}{522 U.S. at 233, 239 n.12}), State~B \textit{must} give Full Faith and Credit to the judgment, notwithstanding State~B's contrary public policy, because the rendering court had jurisdiction over the bound parties and adjudicated a disposition within its authority.

\end{document}
